\chapter{Методы исследования} \label{ch1}

\section{Математическая постановка проблемы} \label{ch1:math}

Пусть задано множество тестовых фреймворков
\[
    F=\{\text{Jest},\text{Vitest},\text{Mocha}\}.
\]
Каждый фреймворк $f \in F$ запускает один и тот же набор модульных тестов $T(D)$ над входным набором данных $D$. Набор данных содержит числовую выборку
\[
    X = (x_1,x_2,\ldots,x_n),
\]
где $n$ --- размер входа. В проекте $D$ может быть импортирован из JSON- или CSV-файла.

Для каждого запуска измеряются три метрики:
\[
    \tau_f(D) \text{ --- время выполнения тестового раннера, мс;}
\]
\[
    m_f(D) \text{ --- пиковое потребление оперативной памяти, МБ;}
\]
\[
    c_f(D) \text{ --- пиковая загрузка CPU, \%.}
\]
Если для каждого фреймворка выполняется $k$ повторов, то наблюдения имеют вид
\[
    r_{f,j}(D)=(\tau_{f,j},m_{f,j},c_{f,j}),\quad j=1,\ldots,k.
\]
Так как отдельный запуск зависит от состояния операционной системы, для итогового сравнения используется медиана:
\[
    \tilde{\tau}_f = median(\tau_{f,1},\ldots,\tau_{f,k}),
\]
\[
    \tilde{m}_f = median(m_{f,1},\ldots,m_{f,k}),
    \quad
    \tilde{c}_f = median(c_{f,1},\ldots,c_{f,k}).
\]

Фреймворк $f_a$ считается быстрее $f_b$ на наборе $D$, если
\[
    \tilde{\tau}_{f_a}(D) < \tilde{\tau}_{f_b}(D).
\]
Аналогичные отношения вводятся для памяти и CPU. В работе не вводится единый ``абсолютный'' рейтинг, потому что в разных проектах важны разные критерии: для локальной разработки чаще важна скорость, для CI --- память и CPU, а для долгосрочной поддержки --- зрелость экосистемы.

\section{Тестируемая функция} \label{ch1:function}

Чтобы сравнить именно стоимость тестового фреймворка, а не сложность предметной области, была выбрана простая, но масштабируемая функция анализа числовой выборки. Для массива $X$ вычисляются:
\begin{enumerate}
    \item количество элементов;
    \item минимум и максимум;
    \item сумма и среднее значение;
    \item медиана;
    \item 90-й процентиль;
    \item число уникальных значений;
    \item количество простых чисел.
\end{enumerate}

Эта функция удобна для эксперимента по двум причинам. Во-первых, ее результат легко проверить на малом эталонном файле. Во-вторых, размер входа можно увеличивать через поле \texttt{size} в JSON-файле, не меняя тесты.

\section{Сравниваемые фреймворки} \label{ch1:frameworks}

Jest используется как комплексный фреймворк модульного тестирования. Он включает собственный раннер, функцию \texttt{expect}, средства mock-тестирования и типичный сценарий запуска через команду \texttt{jest}~\cite{jest-docs}.

Vitest выбран как современный раннер, совместимый по стилю с Jest и рассчитанный на проекты, использующие Vite. В документации Vitest основной режим запуска тестов задается командой \texttt{vitest run}~\cite{vitest-guide,vitest-cli}.

Mocha выбран как минималистичный раннер. В отличие от Jest и Vitest, он не навязывает конкретную библиотеку утверждений. В проекте используется стандартный модуль \texttt{node:assert/strict}, чтобы не добавлять лишнюю зависимость и сравнивать именно раннеры~\cite{mocha-docs,node-assert}.

\noindent
\begingroup
\centering
\small
\begin{longtable}[c]{|p{3.1cm}|p{4.1cm}|p{4.1cm}|p{4.1cm}|}
    \caption{Характеристика сравниваемых фреймворков}
    \label{tab:frameworks}
    \\
    \hline
    Критерий & Jest & Vitest & Mocha \\ \hline
    \endfirsthead
    \hline
    Критерий & Jest & Vitest & Mocha \\ \hline
    \endhead
    \hline
    Тип & Полный фреймворк & Раннер для современной ESM/Vite-среды & Минималистичный раннер \\ \hline
    Утверждения & Встроенный \texttt{expect} & Встроенный \texttt{expect} & Любая библиотека или \texttt{assert} \\ \hline
    Основной запуск & \texttt{jest} & \texttt{vitest run} & \texttt{mocha} \\ \hline
    Ожидаемый профиль & Умеренная память, развитая инфраструктура & Более тяжелая инфраструктура Vite/пула воркеров & Наименьшая инфраструктурная нагрузка \\ \hline
\end{longtable}
\normalsize
\endgroup

\section{Алгоритм измерения} \label{ch1:algorithm}

Псевдокод измерения приведен в алгоритме~\ref{alg:benchmark}. Измерение выполняется отдельным Node.js-скриптом. Для каждого фреймворка создается дочерний процесс, а его PID периодически опрашивается библиотекой \texttt{pidusage}, которая возвращает CPU и память процесса~\cite{pidusage}.

\begin{algorithm}[h]
\caption{Сравнение фреймворков модульного тестирования}
\label{alg:benchmark}
\KwInput{Файл данных $D$, множество фреймворков $F$, число повторов $k$}
\KwOutput{Таблица измерений $R$}
Загрузить $D$ из JSON или CSV\;
Построить числовую выборку $X$\;
$R \leftarrow \varnothing$\;
\ForEach{$f \in F$}{
    \For{$j \leftarrow 1$ \KwTo $k$}{
        Запустить тесты фреймворка $f$ как дочерний процесс\;
        $t_0 \leftarrow$ текущее время\;
        \While{процесс выполняется}{
            считать память процесса по PID\;
            считать загрузку CPU процесса по PID\;
            обновить пиковые значения $m_{max}$ и $c_{max}$\;
        }
        $t_1 \leftarrow$ текущее время\;
        записать $(f,j,t_1-t_0,m_{max},c_{max})$ в $R$\;
    }
}
Сохранить $R$ в JSON и CSV\;
\end{algorithm}

\section{Иллюстративный пример} \label{ch1:example}

Эталонный входной файл \texttt{reference-input.json} содержит массив:
\begin{lstlisting}[breaklines=true]
{
  "name": "reference-json-small",
  "size": 5,
  "values": [9, 1, 5, 3, 7]
}
\end{lstlisting}

Для него результат вычисляется вручную:
\[
    sorted(X)=[1,3,5,7,9],
\]
\[
    count=5,\quad min=1,\quad max=9,\quad sum=25,\quad mean=5,
\]
\[
    median=5,\quad p90=8.2,\quad uniqueCount=5,\quad primeCount=3.
\]
Эти значения записаны в \texttt{reference-output.json}. Команда \texttt{npm run verify} сравнивает вычисленный результат с эталоном.

\section{Пошаговое выполнение метода} \label{ch1:steps}

Пошаговое выполнение метода на эталонных данных:
\begin{enumerate}
    \item программа читает файл \texttt{reference-input.json};
    \item загрузчик определяет формат JSON и извлекает массив \texttt{[9,1,5,3,7]};
    \item модуль анализа сортирует массив и вычисляет статистики;
    \item результат сравнивается с \texttt{reference-output.json};
    \item после успешной проверки запускается бенчмарк;
    \item для Jest, Vitest и Mocha создаются отдельные дочерние процессы;
    \item во время работы процесса сэмплируются RSS-память и CPU;
    \item после завершения каждого процесса фиксируются время, память, CPU и код завершения;
    \item итоговые таблицы сохраняются в каталог \texttt{reports}.
\end{enumerate}
